%
%  chapter6.tex — Conclusion and Future Work
%
\chapter{Conclusion and Future Work}
\label{cha:conclusion}

\section{Summary of Contributions}
\label{sec:summary}

This report presents the design and ongoing implementation of Play4Change, an adaptive learning platform powered by generative AI and gamification mechanics. The project has produced the following verified contributions:

\begin{enumerate}
    \item \textbf{A product concept} addressing a genuine gap: combining AI-generated content, adaptive delivery, peer review, and community engagement in a single platform aligned with the UN Sustainable Development Goals.
    \item \textbf{A cross-platform mobile architecture} using Kotlin Multiplatform, Decompose, and the MVI pattern — sharing all business logic across iOS and Android while preserving native UI on each platform.
    \item \textbf{A type-safe system boundary} via a shared Kotlin module using \texttt{kotlinx.serialization}, eliminating client-server integration bugs at compile time.
    \item \textbf{A railway-oriented error handling model} using Arrow \texttt{Either<AppError, T>}, first built from scratch for comprehension and then migrated to the industry-standard library — demonstrating both understanding and pragmatism.
    \item \textbf{An AI content pipeline} (LangChain4j + Mistral + pgvector RAG) that reduces the educator's authoring burden and self-improves through similarity-based reuse of adaptive branch patterns.
    \item \textbf{A passwordless authentication system} (magic link + OAuth 2.0 with JWKS verification) with security hardening: SHA-256 token storage, atomic CTE claim, token family revocation, and Google audience validation.
    \item \textbf{A peer review assessment mechanism} using three-verdict majority vote for open-ended tasks — zero marginal AI cost, with pedagogical benefit to reviewers through retrieval practice.
    \item \textbf{A complete observability stack} (Micrometer + Prometheus + Grafana) with domain-specific metrics — fully provisioned from source, operational from \texttt{docker compose up}.
    \item \textbf{A documented decision trail} of 16 Architecture Decision Records covering every major engineering choice, including alternatives considered and migration paths to future architectures.
    \item \textbf{A cost-conscious, healthy infrastructure design}: Docker Compose deployment, pgvector over a dedicated vector database, batch AI generation that scales with domains not users, Redis caching with silent fallback, and explicit Kubernetes migration triggers.
\end{enumerate}

\section{Current Limitations}
\label{sec:limitations}

\begin{itemize}
    \item \textbf{No production deployment}: the system runs in Docker Compose. Production VPS provisioning is pending.
    \item \textbf{iOS UI incomplete}: the SwiftUI layer lags behind Jetpack Compose due to Kotlin/Native interoperability complexity.
    \item \textbf{Gamification partially implemented}: streaks and points are operational; badges and leaderboards are pending.
    \item \textbf{Facebook OAuth not hardened}: audience validation via \texttt{/debug\_token} requires \texttt{FACEBOOK\_APP\_ID} and \texttt{FACEBOOK\_APP\_SECRET}, not yet configured (ADR-016, G4 deferred).
    \item \textbf{Two known implementation gaps}: fresh adaptive branch embeddings not stored post-generation; \texttt{runBlocking} inside \texttt{@Transactional} in \texttt{CourseService}. Both are documented in ADR-007 with explicit fix descriptions.
\end{itemize}

\section{Future Work}
\label{sec:future}

\subsection{Kubernetes Deployment}
Migration from Docker Compose to Kubernetes is deferred until one of four explicit triggers is met (ADR-009): resource contention between batch and API serving, multi-institution data residency requirements, zero-downtime deployment SLA, or LLM inference becoming the system bottleneck. The module boundary designed in ADR-006 is already the correct seam for extracting the AI agent into an independently scalable service — zero changes to \texttt{:server} are required.

\subsection{CDN Integration}
Static assets (landing page, topic media, educator-uploaded PDFs served for download) will be fronted by a CDN. The provider decision (Cloudflare, AWS CloudFront, BunnyCDN) has not been made; this is a configuration-level change with no application code impact.

\subsection{Gamification Completion}
The badge award system (achievement triggers, badge definitions, profile display) and leaderboard computation (topic-scoped and global, Redis-cached, daily recomputation) are the next implementation priority.

\subsection{Advanced Analytics and Reporting}
An educator-facing analytics dashboard will show cohort engagement, challenge completion rates, SDG topic coverage, and struggle pattern distributions — enabling educators to identify conceptual gaps and improve topic materials.

\subsection{Richer Adaptive Model}
\begin{itemize}
    \item \textbf{Spaced repetition}: schedule challenge reviews based on the learner's forgetting curve (SM-2 algorithm).
    \item \textbf{Learner interest modeling}: bias content selection toward the SDG tags the learner engages with most.
    \item \textbf{Multimodal challenges}: audio and image-based tasks for learners who respond better to non-textual content.
\end{itemize}

\subsection{Internationalization Extension}
Current support: English and Portuguese. Extending to French and German is straightforward: add \texttt{messages\_fr.properties} on the server and \texttt{values-fr/strings.xml} on the client. Mistral supports content generation in multiple languages.

\subsection{Community and Social Features}
Group challenges, educator communities, and an open topic marketplace are planned for a future phase once the core learning loop reaches production.

\subsection{Apple Sign-In and Additional OAuth Providers}
Adding new OAuth providers requires one new \texttt{OAuthVerifierPort} implementation, one new \texttt{AuthProvider} enum value, and one route registration — zero changes to \texttt{OAuthService}, \texttt{TokenService}, or any domain model (ADR-011).

\section{Conclusion}
\label{sec:conclusion}

Play4Change demonstrates that it is possible to build an educational platform that is simultaneously engaging, architecturally rigorous, cost-sustainable, and secure. The system makes deliberate engineering choices at every layer: type-safe contracts across the client-server boundary, a provider-agnostic AI module that can swap LLMs with a single configuration change, a peer review assessment mechanism that costs nothing at scale while improving reviewer learning, and an observability stack that makes the entire platform's health visible from day one.

The 16 Architecture Decision Records in Appendix~\ref{app:adrs} document not only what was decided but why, what alternatives were considered, and when each decision should be revisited as the system grows. This decision trail is itself a contribution: it transforms the project from a working implementation into a reproducible, extensible system that future engineers can reason about without re-deriving the original constraints.

The foundation is complete. The remaining work — gamification, iOS UI, production deployment, and the richer adaptive model — will not change the architecture; it will validate it.
